\section{Overview of the main results}\label{sec.overview}

Given a commutative superalgebra $A=A^{0}\oplus A^{1}$ over a f\/ield of characteristic 0 and a linear operator $f\colon A\to A$, the hierarchy of higher antibrackets of $f$ is the sequence of operators $\Phi_f^n\colon A^{\odot n}\to A$, $n>0$, def\/ined by the formulas \cite{koszul85,markl2}:
\begin{gather}
\Phi^1_f(a)= f(a),\nonumber\\
\Phi^2_f(a,b)= f(ab)-f(a)b-(-1)^{|a||b|}f(b)a,\nonumber\\
\Phi^3_f(a,b,c)= f(abc)-f(ab)c-(-1)^{|a|(|b|+|c|)}f(bc)a-(-1)^{|b||c|}f(ac)b\nonumber\\
\hphantom{\Phi^3_f(a,b,c)=}{} +f(a)bc+(-1)^{|a||b|}f(b)ac+(-1)^{|c|(|a|+|b|)} f(c)ab ,\nonumber\\
\cdots\cdots\cdots\cdots\cdots\cdots\cdots\cdots\cdots\cdots\cdots\cdots\cdots\cdots\cdots\cdots\cdots\cdots\cdots\cdots\cdots\nonumber\\
\Phi^n_f(a_1,\ldots,a_n)= \sum_{k=1}^n(-1)^{n-k}\sum_{\sigma\in S(k,n-k)}\epsilon(\sigma)
f(a_{\sigma(1)}\cdots a_{\sigma(k)})a_{\sigma(k+1)}\cdots a_{\sigma(n)} ,\label{equ.defiKoszul}
\end{gather}
where $A^{\odot n}$ is the symmetric $n$th power of $A$, where $|a|=0,1$ is the parity of a nonzero homogeneous element $a\in A$, and
$S(k,n-k)$ is the set of shuf\/f\/les of type $(k,n-k)$, i.e., the set of permutations $\sigma$ of $1,\ldots,n$
such that $\sigma(1)<\cdots<\sigma(k)$ and $\sigma(k+1)<\cdots<\sigma(n)$. The Koszul sign
$\epsilon(\sigma)$ is equal to $(-1)^{\alpha}$, where $\alpha$ is the number of pairs $(i,j)$ such that
$i<j$, $\sigma(i)>\sigma(j)$ and $|a_i||a_j|=1$. The same operators can be conveniently def\/ined by the recursive formulas \cite{akman}:
\begin{gather}
\Phi^1_f(a)=f(a),\nonumber\\
\Phi^{n+1}_f(a_1,\ldots,b,c) = \Phi^{n}_f(a_1,\ldots,bc)-\Phi^{n}_f(a_1,\ldots,b)c-(-1)^{|b||c|}\Phi^{n}_f(a_1,\ldots,c)b,\label{equ.defiakman}
\end{gather}
and dif\/fer from the def\/inition given in \cite{alfaro,BDA} by an extra sign factor appearing because in their set-up the operator $f$ acts on the right.
It is immediate to observe that every higher antibracket is a supersymmetric operator, i.e., we have
\begin{gather*} \Phi_f^n(\ldots,a_{i+1},a_i,\ldots)=(-1)^{|a_i| |a_{i+1}|}\Phi_f^n(\ldots,a_{i},a_{i+1},\ldots)\end{gather*}
for every index $0<i<n$ and every sequence $a_1,\ldots,a_n$ of homogeneous elements.
When $A$ has a unit $1\in A^0$, it is well known, and in any case easy to prove, that $\Phi^{n+1}_f=0$
if and only if $f(1)=0$ and $f$ is a dif\/ferential operator of order $\le n$~\cite{koszul85}.

The importance of higher antibrackets relies essentially on the fact that, up to a degree shifting, they satisfy the generalized Jacobi identities \cite{BDA,kravchenko2000}: this means that for every pair of operators
$f$, $g$ we have
\begin{gather*} \Phi^{n}_{[f,g]}=\sum_{i=1}^n\big[\Phi_f^i,\Phi_g^{n+1-i}\big]\end{gather*}
for every $n$, where $[-,-]$ denotes the Nijenhuis--Richardson bracket (see Section~\ref{sec.NR}) on the Lie superalgebra of symmetric operators
\begin{gather*}D(A)=\Hom^*\big(\oplus_{n>0}A^{\odot n},A\big) .\end{gather*}
The Lie superalgebra $D(A)$ contains as special even elements the multiplication operators
\begin{gather*} \mu_n\colon \ A^{\odot n+1}\to A,\qquad \mu_n(a_0,\ldots,a_n)=a_0\cdots a_n ,\qquad n\ge 0 .\end{gather*}
The associated adjoint operators
\begin{gather*} \rho_n\colon \ D(A)\to D(A),\qquad \rho_n(\omega)=[\mu_n,\omega]  \end{gather*}
are even derivations. A tedious but straightforward computation shows that:
\begin{gather*}
\bullet \ \Phi^2_f=-\rho_1(f)=-[\mu_1,f];\\
\bullet \ \Phi^3_f=\dfrac{1}{2}\big(\rho_1^2+\rho_2\big)(f)=\dfrac{1}{2}([\mu_1,[\mu_1,f]]+[\mu_2,f]);\\
\bullet \ \Phi^4_f=-\dfrac{1}{6}\big(\rho_1^3+3\rho_1\rho_2+6\rho_3\big)f=
-\dfrac{1}{6}([\mu_1,[\mu_1,[\mu_1,f]]]+3[\mu_1,[\mu_2,f]]+6[\mu_3,f]);\\
\bullet \ \Phi^5_f=\left(\dfrac{1}{30}\rho_1^4+\dfrac{1}{5}\rho_1^2\rho_2+\rho_1\rho_3+3\rho_4\right)f;\\
\bullet \ \Phi^6_f=-\dfrac{1}{3}\left(\dfrac{1}{90}\rho_1^5+\dfrac{1}{9}\rho_1^3\rho_2+\rho_1^2\rho_3+7\rho_1\rho_4+28\rho_5\right)f.
\end{gather*}
The above expressions for $\Phi^2_f$ and $\Phi^3_f$ were also pointed out and conveniently used in~\cite{FMpoisson} in the framework of Poisson geometry.

It is natural to ask whether similar formulas hold for every $n$: more precisely, we ask for a~sequence of noncommutative polynomials $P_n(\rho_1,\ldots,\rho_n)$ with rational coef\/f\/icients such that, for every~$f$, we have
\begin{gather}\label{equ.polinomigeneratori}
 \Phi_f^{n+1}=P_n(\rho_1,\ldots,\rho_n)f .
 \end{gather}
